\section[Weak nonlinear learning in q20]{When transport training learns scale but little nonlinear shape}
\label{sec:bf-neutra-q20-repair}

The September 2026 q20 investigation asks why reverse-KL training can produce
an almost affine map while substantial nonlinear posterior geometry remains.
The target is the four-parameter, thirty-observation UKF approximate posterior
of the SSL--LSTM example. The ``q20'' designation refers to its latent-state
configuration, not a twenty-dimensional parameter vector. All statements in
this section concern that specified approximate posterior; accurate sampling
from it would not establish agreement with the exact state-space posterior.
The earlier three-stage examples in this chapter and the two/four-stage q20
recipes are distinct experiments.

\subsection{What failed, and what has already been checked}

Let $z\sim r=\calN(0,I_4)$, let $\theta=T_\phi(z)$, and define the
unnormalized transformed log density and its Gaussian-score residual as
\begin{equation}
 \ell_Z(z)=\ell_\Theta(T_\phi(z))+\log|\det J_T(z)|,
 \qquad g(z)=\nabla_z\ell_Z(z)+z.
 \label{eq:bf-q20-residual}
\end{equation}
An exact Gaussianizing map has $g=0$. A useful HMC map need not make this
identity exact, so the residual explains geometry and can nominate a repair;
posterior estimation still requires the sampling checks described elsewhere
in this chapter.

On the saved bank of 1,000 base-normal points, affine regressions explain
$99.866$--$99.999\%$ of the coordinate variation of the depth-four map. Of
the observed squared residual $\sum_n\|g(z^{(n)})\|^2$, $89.02\%$ lies in
latent coordinate 2. A constant-plus-linear regression on all four latent
coordinates explains only $0.573\%$ of that component's residual energy;
its root-mean-square magnitude is $4.2638$. That latent direction principally
controls the physical observation weight, approximately
$\theta_3=0.2743+0.4958z_2$ on this bank. These descriptive observations
localize a failure to learn nonlinear shape. They are not a bound on derivatives
of the fitted map, proof of separability, or evidence that the target has
multiple modes.

The corresponding source audit checked the loss
\begin{equation}
 -\E_r[\ell_\Theta(T_\phi(z))+\log|\det J_T(z)|]
\end{equation}
and verified that its first
parameter derivative includes both the supplied analytic target score and
the log determinant. The previously checked finite differences do not support
a missing-score explanation. The reported loss near 43 is not a KL divergence
of 43: the additive terms $\E_r\log r+\log Z$ are not included in that number.

Clipping alone also cannot explain the remaining failure. The continuation
control clipped $99.02\%$ of its updates, whereas the clipping repair clipped
$0.68\%$ and the depth-four continuation clipped none. Each principal repair
added 1,024 updates with batch 32. The added layers consequently have only
1,024 updates, and the existing layers have 2,048 lifetime updates. The late
depth-four losses are flat and noisy; improvement from parent to endpoint does
not establish continuing improvement at the endpoint or optimizer convergence.

These results and their frozen-source identities are preserved in the
\href{../plans/bayesfilter-q20-training-math-audit-results-2026-09-23.md}{September 23 mathematical audit}.
They motivate three distinct mechanisms: release overall scale from a
saturating parameterization, reduce noise in the same reverse-KL gradient,
and expose direct scalar shape changes. The observations do not yet determine
which mechanism matters most.

\subsection{Why reverse-KL stationarity need not remove the score residual}

The objective averages log density, whereas $g$ measures its derivative.
Even an exact optimum within an affine family can retain nonlinear score
error. Consider $p(x)\propto\exp(-x^4/12)$ and $x=m+\sigma z$. The
nonconstant reverse-KL objective is
\begin{equation}
 L(m,\sigma)=\frac{m^4+6m^2\sigma^2+3\sigma^4}{12}-\log\sigma.
\end{equation}
For fixed positive $\sigma$, its minimum in $m$ is zero. At $m=0$,
$\partial_\sigma L=\sigma^3-\sigma^{-1}$, so the optimum is $\sigma=1$.
Nevertheless $g(z)=z-z^3/3$, and the Gaussian moments give
\begin{equation}
 \E[g]=\E[zg]=0,\qquad
 \E[g^2]=1-\tfrac23(3)+\tfrac19(15)=\tfrac23.
\end{equation}
Thus affine directions can be stationary while nonlinear error remains. This
is an analytic example, not a fitted description of the q20 posterior.

Adding $\E\|g\|^2$ to the objective would change the optimization problem.
Its parameter derivative generally requires derivatives of the target score,
which the existing stopped-score bridge does not supply. The following
path-gradient mechanism instead preserves reverse KL and needs only the
first target score.

\subsection{Mechanism one: separate global scale from conditional deformation}

The q20 conditional log scale is $s(a)=2\tanh(a/2)$. At the observed mean
scale near $-1.9936$,
\begin{equation}
 \frac{ds}{da}=1-\left(\frac{s}{2}\right)^2\simeq0.0064.
\end{equation}
This is a measured local attenuation. It does not prove the composed map
cannot contract: other stages can compensate, and Adam's update is not
proportional to this derivative alone. Still, the fixed outer scale of four
requires contraction somewhere, so making every global scale change pass
through the cap is an avoidable optimization constraint.

One parameterization supported by the original NeuTra author implementation is
\begin{equation}
 s_j(x_{<j})=b_j+c\tanh\!\left(\frac{h_j(x_{<j})}{c}\right),
 \qquad \frac{\partial s_j}{\partial b_j}=1.
 \label{eq:bf-q20-free-scale}
\end{equation}
The bounded term controls input-dependent variation and the free bias controls
global contraction. A final cap on their sum would restore the obstruction.
The inspected \path{neutra/utils.py} implementation applies its optional
\texttt{log\_scale\_clip\_pre} before adding \texttt{log\_scale\_bias};
this is an available source implementation, not a claim that every published
experiment used that option. Andrade's stable-flow construction similarly
separates conditional clamps from a final trainable affine transformation
\citep[Sections 3--5]{andrade2024stableflows}.

The smallest map-preserving implementation is a separate affine correction
\begin{equation}
 A_\eta(x)=m+\diag(e^b)x,\qquad
 \log|\det J_A|=\sum_j b_j,
 \qquad m=b=0\ \text{at initialization}.
 \label{eq:bf-q20-affine-mechanism}
\end{equation}
Then $A_\eta\circ T_0=T_0$ initially. This adds unconstrained global scale
without rewriting the old checkpoint or optimizer slots. Its finite-state
checks must reject overflow rather than clip the posterior or omit offending
draws. It can improve scale conditioning; by itself it does not supply the
missing nonlinear scalar shape.

\subsection{Mechanism two: an exact alternative reverse-KL gradient}

Write $s_q(\theta)=\nabla_\theta\log q_\phi(\theta)$ and
$s_\pi(\theta)=\nabla_\theta\ell_\Theta(\theta)$. With a differentiable
full-support map and integrability sufficient to interchange differentiation
and integration, the chain rule gives
\begin{align}
 \nabla_\phi\operatorname{KL}(q_\phi\Vert\pi_\Theta)
 &=\E_r\big[(\partial_\phi T)^\top(s_q-s_\pi)\big]
   +\E_q\big[\partial_\phi\log q_\phi(\theta)|_\theta\big],\\
 \E_q\big[\partial_\phi\log q_\phi(\theta)|_\theta\big]
 &=\partial_\phi\int q_\phi(\theta)\,d\theta=0.
 \label{eq:bf-q20-path-gradient}
\end{align}
The path estimator retains the first expectation. At $q=\pi$ it vanishes
for each draw, whereas the ordinary estimator can remain noisy. Removing a
zero-mean term does not guarantee lower variance away from the optimum,
because covariance between terms also matters
\citep[Section 2]{roeder2017sticking}. A finite minibatch therefore need not
give equal standard and path gradients. Their expectations, rather than their
pointwise values, must be checked against the same objective.

The proposal score can be obtained without numerical inversion. Differentiate
the change-of-variables identity at a known base draw:
\begin{align}
 \log q_\phi(T_\phi(z))&=\log r(z)-\log|\det J_T(z)|,\\
 J_T(z)^\top s_q(T_\phi(z))&=-z-\nabla_z\log|\det J_T(z)|.
 \label{eq:bf-q20-forward-proposal-score}
\end{align}
Solving this transposed-Jacobian system supplies $s_q$. Applied one layer at
a time it is the score recursion of
\citet[Proposition 3.2 and Appendix B.1]{vaitl2024fastpath}. Their optimized
coupling-flow timing results do not establish the cost of our autoregressive
implementation. The first optional implementation uses a bounded small dense
Jacobian assembled by coordinate-wise vector-Jacobian products. That loop is
over the four parameter coordinates, not training samples; all sample rows
remain batched. It uses no TensorFlow pfor and makes no coupling-flow speed claim.

After computing $v=s_q-s_\pi$, hold $v$ fixed and differentiate
$B^{-1}\sum_n T_\phi(z^{(n)})^\top v^{(n)}$ with respect to $\phi$.
This scalar carries the intended gradient but is not the reverse-KL loss;
the actual $-\ell_\Theta-\log|\det J_T|$ must be reported separately.
Stopping the complete forward log determinant would discard necessary terms.
The inverse-density construction of \citet{vaitl2022path} and the forward
recursion above are two ways to compute the required proposal score correctly.

Combining \eqref{eq:bf-q20-residual} and
\eqref{eq:bf-q20-forward-proposal-score} yields
\begin{equation}
 J_T^\top(s_q-s_\pi)=-g,\qquad
 \nabla_\phi\operatorname{KL}
 =-\E_r\big[(J_T^{-1}\partial_\phi T)^\top g\big].
 \label{eq:bf-q20-residual-projection}
\end{equation}
If available parameter variations have little projection onto $g$, the expected
gradient can be small despite large nonlinear error. Variance reduction helps
resolve an existing direction; it does not create a missing one.

\subsection{Mechanism three: direct scalar deformation with an inexpensive log determinant}

\citet[Sections 2--3 and 6.1.2]{huang2018naf} identify poor fits and local
minima with conditional affine flows and replace the scalar affine function
by a monotone neural transformation. For a deep sigmoid flow (DSF), let
\begin{equation}
 \begin{aligned}
 S(x)=\sum_{k=1}^{K}w_k\sigma(a_kx+b_k),\quad
 &\tau(x)=\log S(x)-\log(1-S(x)),\\
 a_k>0,\quad w_k>0,\quad\sum_kw_k=1.
 \end{aligned}
 \label{eq:bf-q20-dsf}
\end{equation}
Here $\sigma$ is logistic. Positivity makes the derivative strictly positive;
the limits are $\tau(x)\to\pm\infty$ as $x\to\pm\infty$. For $K=1$,
$\tau(x)=a_1x+b_1$, giving an exact affine control.

Let $p_k=\sigma(a_kx+b_k)$ and
$N=\sum_kw_ka_kp_k(1-p_k)$. Direct differentiation gives
\begin{equation}
 \tau'(x)=\frac{N}{S(1-S)},\qquad
 \log\tau'(x)=\log N-\log S-\log(1-S).
 \label{eq:bf-q20-dsf-logdet}
\end{equation}
Thus there is a fast analytic determinant formula. When the scalar parameters
depend only on preceding coordinates, the full autoregressive Jacobian remains
triangular and $\log|\det J|=\sum_j\log\tau'_j$. A masked conditioner
produces all coordinate parameters in one pass. A single DSF scalar layer costs
$O(BdK)$ beyond the conditioner; a deep dense scalar network also incurs its
hidden matrix-product costs. These operation counts are not measured GPU timings.

Directly forming $1-S$ loses precision in the tail. With
$u_k=a_kx+b_k$, calculate instead
\begin{align}
 A&=\operatorname{LSE}_k(\log w_k+\log\sigma(u_k)),\\
 B_0&=\operatorname{LSE}_k(\log w_k+\log\sigma(-u_k)),\\
 C&=\operatorname{LSE}_k(\log w_k+\log a_k+
                         \log\sigma(u_k)+\log\sigma(-u_k)),\\
 \tau&=A-B_0,\qquad \log\tau'=C-A-B_0.
 \label{eq:bf-q20-dsf-stable}
\end{align}
Here $\operatorname{LSE}(v)=\log\sum_k\exp v_k$ is evaluated stably.
This algebra is consistent with the log-domain approach in the paper's
Appendix C. The inspected author code additionally shrinks $S$ to
$(1-\delta)S+\delta/2$ for numerical stability. That yields finite output
limits and differs from the full-range map required here; the proposed
implementation uses \eqref{eq:bf-q20-dsf-stable} without that shrinkage.

HMC and geometry diagnostics need the derivative of the log determinant too.
Let $\rho_k=w_ka_kp_k(1-p_k)/N$. Then
\begin{equation}
 \frac{\tau''}{\tau'}
 =\sum_k\rho_ka_k(1-2p_k)-\exp(C-A)+\exp(C-B_0).
 \label{eq:bf-q20-dsf-force}
\end{equation}
The first term is $N'/N$, and the last two differentiate $-\log S$ and
$-\log(1-S)$. This remains a batched scalar reduction. For a conditional
map, derivatives through the conditioner also contribute; the scalar formula
alone is not the entire multidimensional score.

NAF's significant cost distinction is its inverse. Given $y=\tau(x)$, each
component logit crosses $y$ at $r_k=(y-b_k)/a_k$. Therefore
\begin{equation}
 \min_k r_k\ \leq\ \tau^{-1}(y)\ \leq\ \max_k r_k.
 \label{eq:bf-q20-dsf-bracket}
\end{equation}
At the lower bound every sigmoid is at most $\sigma(y)$; at the upper bound
every sigmoid is at least $\sigma(y)$, proving the bracket. A bounded bisection
can solve inside it, with explicit input/output tolerances, an iteration cap,
and a failure result. Training in the forward reverse-KL direction and ordinary
latent HMC steps do not require this inverse on every evaluation. Physical-space
proposal densities, forward-KL training, and ensemble cross densities can do so.
Fast forward determinants consequently do not imply a cheap complete ensemble
pipeline. The first implementation does not differentiate through bisection or
offer inverse-based training.

There is also an initialization trap. At $a_k=1,b_k=0$ for all components,
with $w=\operatorname{softmax}(\eta)$,
\begin{equation}
 \partial_{a_k}\tau=w_kx,\qquad
 \partial_{b_k}\tau=w_k,\qquad
 \partial_{\eta_k}\tau=0.
\end{equation}
The scalar parameter directions are initially only affine. Small distinct
components expose nonlinear directions, but their initial map displacement
must be measured. Zero output initialization is not generally a defect:
successful NAF and Glow recipes start near identity
\citep[Appendix E]{huang2018naf}\citep[Section 3.3]{kingma2018glow}.
The question is whether the chosen features and parameter directions can
respond usefully on this target.

The smallest shape experiment freezes the current map and uses
$T_\psi=T_0\circ D_\psi$, where
$D_\psi(z)=(z_1,\tau_\psi(z_2),z_3,z_4)$. It includes
$\log\tau'_\psi(z_2)$ in the loss and keeps the posterior unchanged.
The observed concentration of error motivates coordinate 2, but does not
guarantee a marginal correction will suffice. Residual dependence on other
coordinates would motivate a conditional transform next.

Explicit inverses are available for rational-quadratic splines, along with
inexpensive determinants \citep[Section 3]{durkan2019splines}. Their usual construction is
$C^1$ but generally not $C^2$. Since the scalar transformed force contains
$\tau''/\tau'$, knot behavior requires a separate HMC energy and derivative
check. They remain an alternative, not an automatically qualified replacement.

\subsection{How the three mechanisms will be distinguished}

First test the arithmetic on analytic targets: affine identities, Gaussian
pointwise cancellation of the path gradient, equality of expected standard and
path gradients under converged quadrature, finite differences of scalar values
and log determinants, and inverse roundtrips including tails. Checkpoint
reconstruction must preserve all mechanism parameters and never imply HMC
admission. Tiny CPU fixtures with hidden GPUs answer these engineering
questions only. A few successful optimizer steps are not training evidence
for the q20 posterior.

The target experiment then starts from the same saved endpoint for isolated
scale, gradient-estimator, and scalar-shape arms, with unchanged reverse KL.
An affine-only scalar control distinguishes shape from location and scale.
Pair the latent streams where meaningful and charge actual worker time; equal
target-row counts need not imply equal optimizer steps or equal runtime.
Keep true losses, per-block gradient/update measurements, scale derivatives,
inverse work where used, and checkpoints. A failed candidate triggers the
corresponding next repair; it does not alone reject NeuTra.

Every numerical choice needs a local basis. The coordinate follows the saved
residual decomposition. The minimum nonlinear component count follows from
$K=1$ being affine, but sufficient capacity is uncalibrated. Positive slopes
and weights follow from invertibility. Initial component spread is judged by
map displacement and nonlinear response. Batch, learning rate, continuation
length, replication, and any revised conditional cap must be priced and
calibrated on the actual target. Published values such as the original NeuTra
batch 4,096 and 5,000 updates are experimental settings, not q20 defaults.

Fresh heldout objective estimates and the standard 1,000-point score report
can support nomination. They cannot establish mode coverage, convergence, or
superiority among stochastic candidates without the corresponding evidence.
A new scalar family also requires supported frozen-map reconstruction,
transformed-score checks, and fresh public-tuner qualification before HMC.
The implementation and experiment sequence is recorded in the
\href{../plans/bayesfilter-q20-three-mechanisms-plan-2026-09-23.md}{three-mechanism repair plan}.
The strongest competing explanation remains that the current IAF family could
learn adequate nonlinear shape under a better calibrated continuation. The
three tests are designed to distinguish that explanation from limited effective
parameter directions, rather than assume an architecture replacement is necessary.

\subsection{A shared implementation of the proposed mechanisms}

The September 24 implementation makes the affine autoregressive map, its
inverse, and its score derivatives common to ordinary, weighted, legacy,
and frozen-map consumers. Historical checkpoints retain their original scale
conventions. A separate source-based IAF configuration implements the NeuTra
author code's conditional cap with a free bias at every layer, together with
the paper's three-stage, two-hidden-layer ELU architecture. The fixed cap is
not a bound on the free global scale. The source initializer uses fan-in
variance scaling with factor $0.02$, as selected by the author's
\texttt{MakeIAFBijectorFn}; neither that factor nor the paper's training
schedule establishes a calibrated q20 protocol.

The new DSF configuration is a full conditional neural autoregressive flow:
the earlier coordinates determine each coordinate's positive slopes, shifts,
and mixture weights. Its source profile retains the author's cMADE
conditioner, final conditioner weight normalization, and shared output
projection, specialized to a constant context for an unconditional target.
The log-domain transformer preserves the paper's full real support. It omits
the author implementation's output shrinkage, which would bound that support.
This numerical distinction and the fixed random-number conventions are
documented explicitly in the
\href{../reference/neutra-implementation.md}{implementation reference}.

Layerwise proposal-score propagation is connected to the Adam training
update. It uses Vaitl's general score recursion and autoregressive triangular
solves, not the coupling-specific linear-cost formula. Inverse-density
gradients use implicit differentiation of the same conditional map.
Training handoff includes the standard 1,000-point residual report; a finite
report alone establishes neither adequate whitening nor posterior convergence.
These changes close implementation gaps. Their numerical regression and
GPU checks are engineering evidence; a freshly calibrated q20 training run
and downstream sampler validation are still required.
